\title{The Era of 1-bit LLMs: All Large Language Models are in 1.58 Bits}

\begin{abstract}
Advancements in research, notably BitNet~\cite{bitnet}, are heralding the onset of a new epoch for 1-bit Large Language Models (LLMs). In this study, we present a 1-bit LLM variant referred to as \textbf{\text{BitNet b1.58}}, characterized by all parameters (weights) being ternary, taking values from \{-1, 0, 1\}. This model achieves perplexity and downstream task performance on par with full-precision Transformer LLMs (FP16 or BF16) of identical size and exposed to the same number of training tokens, while offering major gains in latency, memory footprint, throughput, and energy efficiency. Importantly, the 1.58-bit LLM establishes a novel scaling law and framework for training future LLMs that maintain high performance while optimizing resource usage. In addition, it enables an innovative computational approach and lays the groundwork for hardware tailored to the unique demands of 1-bit LLMs.
\end{abstract}

\section{The Era of 1-bit LLMs}

The expansion in scale and performance of Large Language Models (LLMs) has been a defining trend in artificial intelligence in recent years. While these models deliver exceptional results across an array of natural language processing benchmarks, their continually growing size introduces considerable hurdles for practical deployment, along with concerns about sustainability due to substantial energy requirements.

A key strategy to mitigate these concerns has involved applying post-training quantization techniques, which convert trained models into lower bitwidth formats for inference~\cite{smoothquant, gptq, quip, quip_sharp}. By reducing both weight and activation precision, such methods sharply decrease the computational and memory burdens of LLMs. The progression has evolved from 16-bit precision toward ever smaller bit-widths, including various 4-bit model implementations~\cite{gptq, awq}. Despite their prevalence in industrial LLM deployments, post-training quantization methods remain inherently limited in performance relative to alternative solutions.

Recent advancements in 1-bit neural architectures, exemplified by BitNet~\cite{bitnet}, have emerged as a viable strategy to reduce the operational costs associated with large language models (LLMs), while preserving high levels of performance. Conventional LLMs typically utilize 16-bit floating-point representations (such as FP16 or BF16), and most of their computational load is dominated by matrix multiplications. Consequently, the high computational expenses predominantly arise from floating-point addition and multiplication. Conversely, BitNet's design limits matrix multiplication to integer addition, resulting in substantial energy savings for LLM operations. Since many modern chips are constrained by power as a fundamental bottleneck for compute throughput, these energy reductions can often translate into increased computational speed.

Beyond raw computation, transferring model parameters from external DRAM into on-chip memory like SRAM during inference also incurs significant costs. Expanding SRAM to alleviate this issue and increase throughput is possible, but comes with much higher expense compared to DRAM. 1-bit LLMs markedly reduce memory requirements—both in capacity and bandwidth—relative to full-precision models. This decreased footprint can considerably lower the expense and latency of loading model weights from DRAM, thereby facilitating more rapid and cost-efficient inference.

Within this context, we propose a substantial evolution of 1-bit LLMs, named \textbf{\text{BitNet b1.58}}, where each model parameter is ternary and can assume values from the set \{-1, 0, 1\}. This extension of the original BitNet framework—by introducing 0 as a possible value—yields an effective 1.58 bits per parameter. \text{BitNet b1.58} upholds all the fundamental advantages of the earlier BitNet, such as its shift to computation that nearly eliminates multiplication in matrix operations, facilitating highly optimized implementations. It maintains equivalent energy efficiency compared to the original 1-bit BitNet, and further surpasses FP16-based LLMs in memory usage, throughput, and latency. Moreover, \text{BitNet b1.58} brings two noteworthy enhancements: its ternary weights enhance the model’s expressive power through direct support for feature selection (enabled by the zero value), which can notably enhance the capabilities of 1-bit LLMs. Our empirical evaluations further demonstrate that \text{BitNet b1.58}, beginning at model sizes of 3B, can achieve parity with full-precision FP16 baselines on both perplexity and downstream tasks, given identical settings in terms of scale and training data.

\section{BitNet b1.58}

\text{BitNet b1.58} extends the BitNet family, employing a Transformer structure where the standard \emph{nn.Linear} operations are substituted by \emph{BitLinear} modules. This model is initialized without any pretraining, using 1.58-bit quantized weights and 8-bit quantized activations. In comparison to the initial BitNet, several refinements are introduced and outlined below.

\paragraph{Quantization Function.} To enforce the constraint that weights take values only in \{-1, 0, +1\}, an \emph{absmean} quantization approach is utilized. This involves rescaling the entire weight matrix by its mean absolute value, followed by mapping each entry to the closest integer in \{-1, 0, +1\}:

\begin{equation}
    \widetilde{W} = \mathrm{RoundClip}(\frac{W}{\gamma + \epsilon}, -1, 1),
\end{equation}

\begin{equation}
    \mathrm{RoundClip}(x, a, b) = \max(a, \min(b, \mathrm{round}(x))),
\end{equation}

\begin{equation}
    \gamma = \frac{1}{nm}\sum_{ij} |W_{ij}|.
\end{equation}

The method used for quantizing activations mirrors the approach from BitNet, with the exception that we omit the scaling of activations to $[0, Q_b]$ before non-linearities. Instead, the activation values are adjusted to $[-Q_b, Q_b]$ for each token individually, eliminating the use of zero-point quantization. This revision streamlines implementation and facilitates system-level optimization, without noticeably impacting model performance in our evaluations.

\paragraph{LLaMA-alike Components.}

LLaMA~\cite{llama, llama2} has emerged as the standard architecture for open large language models in the open-source ecosystem. Aligning with this trend, \text{BitNet b1.58} incorporates components analogous to those in LLaMA. This includes adopting RMSNorm~\cite{rmsnorm}, SwiGLU~\cite{swiglu}, rotary positional encodings~\cite{rope}, and discarding all bias terms. These design decisions enable seamless compatibility with leading open-source libraries such as Huggingface, vLLM~\cite{vllm}, and llama.cpp\footnote{\url{https://github.com/ggerganov/llama.cpp}}, thereby ensuring straightforward adoption and integration.

\section{Results}

We evaluated the performance of \text{BitNet b1.58} alongside our re-implemented FP16 \text{LLaMA LLM} across different parameter scales. Both models were pre-trained from scratch on the RedPajama dataset~\cite{redpajama}, utilizing 100 billion tokens to maintain consistent experimental conditions. For assessment, we measured zero-shot results across an array of language understanding benchmarks, such as ARC-Easy~\cite{arc}, ARC-Challenge~\cite{arc}, Hellaswag~\cite{hellaswag}, Winogrande~\cite{winoGrande}, PIQA~\cite{piqa}, OpenbookQA~\cite{openbookqa}, and BoolQ~\cite{boolq}. Additionally, validation perplexity scores were provided on the WikiText2~\cite{wikitext2} and C4~\cite{c4} corpora.

We also recorded GPU memory usage and inference latency for both \text{LLaMA LLM} and \text{BitNet b1.58}, leveraging the FasterTransformer\footnote{\url{https://github.com/NVIDIA/FasterTransformer}} framework, which is designed for high-performance LLM inference on GPU hardware. For \text{BitNet b1.58}, the Ladder~\cite{ladder} 2-bit kernel was integrated. The primary inference efficiency metric was time taken per output token.

As detailed in Table~\ref{tab:ppl}, we present the perplexity and associated computational costs for both \text{BitNet b1.58} and \text{LLaMA LLM}. The findings indicate that from the 3B parameter scale onward, \text{BitNet b1.58} achieves comparable perplexity to its full-precision \text{LLaMA LLM} counterpart, while offering 2.71$\times$ greater speed and using 3.55$\times$ less GPU memory. Specifically, at the 3.9B model size, \text{BitNet b1.58} surpasses \text{LLaMA LLM} 3B in performance, reduces memory consumption by 3.32$\times$, and is 2.4$\times$ faster.

Table~\ref{tab:zero-shot} contains the zero-shot evaluation metrics on the aforementioned downstream tasks. We adopted the evaluation pipeline from \emph{lm-evaluation-harness}\footnote{\url{https://github.com/EleutherAI/lm-evaluation-harness}}. The results demonstrate that as the model scale increases, the performance disparity between \text{BitNet b1.58} and \text{LLaMA LLM} diminishes. Notably, starting at 3B parameters, \text{BitNet b1.58} attains parity with the full-precision baseline on these tasks. Mirroring the perplexity trends, the downstream results confirm that \text{BitNet b1.58} 3.9B consistently outperforms \text{LLaMA LLM} 3B, while incurring lower memory and latency overheads. These observations underscore that \text{BitNet b1.58} achieves a Pareto improvement over leading LLMs.

\paragraph{Memory and Latency}

We expanded our experiments to include larger models at the 7B, 13B, and 70B parameter scales, assessing the associated computational costs. As depicted in Figure~\ref{fig:latency-memory-energy}, both latency and memory usage exhibit clear scaling patterns, with acceleration becoming more pronounced at higher model sizes. Notably, the \text{BitNet b1.58} 70B variant achieves a speed-up of 4.1$\times$ over the baseline \text{LLaMA LLM}. This improvement stems from the increasing computational burden of \emph{nn.Linear} as model dimensionality grows. Memory demands display a parallel pattern, given that embeddings are kept in full precision and comprise a diminishing share of total memory usage as models increase in size. All latency and memory metrics were gathered using a 2-bit kernel implementation, indicating that further efficiency gains are possible through additional optimization.

\paragraph{Energy}

To assess efficiency at the hardware level, we estimated the energy required for arithmetic operations in both \text{BitNet b1.58} and \text{LLaMA LLM}. Our analysis centers on matrix multiplications, which dominate the computational workload in large language models. Figure~\ref{fig:energy} details the distribution of energy expenditures: in \text{BitNet b1.58}, most operations are INT8 additions, whereas \text{LLaMA LLM} relies on both FP16 additions and multiplications. Based on the energy model described in~\cite{energycost, pokebnn}, the matrix multiplication operations in \text{BitNet b1.58} consume 71.4 times less energy than those in the FP16-based system on 7nm process technology. We also calculated the end-to-end energy usage for inputs of 512 tokens, finding that energy efficiency advantages of \text{BitNet b1.58} relative to the FP16 \text{LLaMA LLM} baseline become even more pronounced with increasing model scale. This effect arises because the relative share of computations devoted to \emph{nn.Linear} increases in larger models, while contributions from other subsystems shrink proportionally.

\paragraph{Throughput}

We evaluated the throughput performance of \text{BitNet b1.58} in comparison with the 70B parameter version of \text{LLaMA LLM}, conducting experiments across two 80GB A100 GPUs by leveraging pipeline parallelism~\cite{gpipe} to accommodate the 70B parameter \text{LLaMA LLM} model. For both models, the batch size was increased incrementally until the GPU memory constraint was reached, with a fixed sequence length of 512 tokens. As reported in Table~\ref{tab:throughput}, \text{BitNet b1.58} 70B is capable of supporting batch sizes up to 11 times greater than those of \text{LLaMA LLM}, leading to an 8.9-fold improvement in overall throughput.

\textbf{\text{BitNet b1.58} introduces a novel scaling paradigm relating model size, inference efficiency, and performance.} Drawing from Figures~\ref{fig:latency-memory-energy} and \ref{fig:energy}, equivalencies can be established between 1.58-bit models and conventional 16-bit models with respect to efficiency:

\begin{itemize}
    \item BitNet b1.58 with 13B parameters surpasses a 3B FP16 LLM in latency, memory footprint, and energy demands.
    \item The 30B variant of BitNet b1.58 outperforms the 7B FP16 LLM across the same efficiency metrics.
    \item At 70B parameters, BitNet b1.58 demonstrates greater efficiency than the 13B FP16 LLM regarding latency, memory usage, and energy consumption.
\end{itemize}

\paragraph{Training with 2T Tokens}

The volume of training tokens plays a pivotal role in large language model performance. To assess the token scalability of \text{BitNet b1.58}, we trained a version of this model on a 2T-token dataset, employing the StableLM-3B~\cite{StableLM-3B-4E1T} data processing strategy, recognized as state-of-the-art for open-source 3B models. Evaluation was conducted using a benchmark suite comprising Winogrande~\cite{winoGrande}, PIQA~\cite{piqa}, SciQ~\cite{sciq}, LAMBADA~\cite{lambada}, and ARC-easy~\cite{arc}. Zero-shot accuracy results are presented in Table~\ref{tab:2t}. For metrics involving both accuracy and normalized accuracy, their average is reported. The outcomes for StableLM 3B at 2T tokens are drawn from its respective technical documentation. Our experimental results indicate that \text{BitNet b1.58} consistently achieves superior performance on all assessed tasks, highlighting the robust generalization capability of 1.58-bit LLMs.

\section{Discussion and Future Work}

\textbf{1-bit Mixture-of-Experts (MoE) LLMs}

Mixture-of-Experts (MoE) architectures offer a cost-efficient solution for scaling LLMs, primarily by decreasing computational FLOPs. However, their adoption remains limited by substantial memory demands and significant inter-chip communication costs. Employing 1.58-bit LLMs can effectively mitigate these issues. Reduced memory usage lowers the hardware requirements needed for MoE deployment. Additionally, transferring activations across devices is notably simplified, and the problem could be eliminated altogether if the model fits within the confines of a single chip.

\textbf{Native Support of Long Sequence in LLMs}

As LLM applications increasingly require processing of extended sequences, efficient long-sequence inference has become a crucial challenge, especially due to the memory overhead from KV caches. The design of \text{BitNet b1.58} marks substantial progress in addressing this issue, by compressing activations from 16-bit to 8-bit, effectively doubling possible context length with unchanged resources. There exists further potential to losslessly compress to 4 bits or even less for 1.58-bit LLMs—a direction we propose for future investigation.

\textbf{LLMs on Edge and Mobile}

Leveraging 1.58-bit LLMs could markedly enhance language model performance on edge and mobile hardware, which frequently faces strict constraints in terms of memory capacity and computational throughput. Owing to their low memory and energy requirements, 1.58-bit LLMs are better suited to these platforms, unlocking applications previously unattainable on such devices. This advance stands to significantly expand the capabilities and utility of LLMs at the edge and on mobile. Furthermore, 1.58-bit LLMs are optimized for CPUs, the predominant processors in these environments, enabling efficient execution of models like \text{BitNet b1.58} and further boosting their performance.

\textbf{New Hardware for 1-bit LLMs}

Recent developments, such as those by Groq\footnote{\url{https://groq.com/}}, highlight the promise of designing dedicated hardware—such as LPUs—tailored for LLM acceleration. Building on this trajectory, we propose the exploration and creation of novel hardware and systems specifically tailored for the computational characteristics of 1-bit LLMs, taking advantage of the unique paradigm introduced in BitNet~\cite{bitnet}.

\end{document}